\documentclass[11pt]{article}

% --- packages ---
\usepackage[margin=1in]{geometry}
\usepackage{amsmath,amssymb,amsfonts}
\usepackage{booktabs}
\usepackage{graphicx}
\usepackage{hyperref}
\usepackage{xcolor}
\usepackage{multirow}
\usepackage{float}
\usepackage[numbers]{natbib}
\usepackage{algorithm}
\usepackage{algorithmic}
% \usepackage{microtype}  % disabled: MiKTeX font expansion issue
\usepackage{bookmark}

\hypersetup{
  colorlinks=true,
  linkcolor=blue!60!black,
  citecolor=blue!60!black,
  urlcolor=blue!60!black
}

% --- metadata ---
\title{%
  La Matrice Ponte Addestrabile:\\
  Fingerprinting del Task e Composizione di Adattatori\\
  tramite Accoppiamento Strutturato nell'Adattamento a Basso Rango%
}

\author{%
  Timothy Paul Bielec%
  \thanks{Autore corrispondente. \texttt{timothy@promptcrafted.com}}%
  \\
  Promptcrafted LLC%
}

\date{Marzo 2026}

\begin{document}
\maketitle

\begin{abstract}
Inseriamo una matrice di accoppiamento addestrabile $6 \times 6$---la
\emph{matrice ponte}---tra le proiezioni discendente e ascendente di un
adattatore LoRA, aggiungendo 36 parametri per ogni livello adattato
($0{,}02\%$ di sovraccarico). La matrice ponte \`e inizializzata come
matrice identit\`a, recuperando esattamente il LoRA standard. Attraverso
sette esperimenti su Qwen2.5 (1,5B e 7B), stabiliamo quattro risultati:
(1)~le matrici ponte codificano l'identit\`a del task---un SVM
leave-one-out su 28 matrici ponte delle proiezioni query (1.008
parametri) classifica il tipo di task con un'accuratezza dell'82,1\%
(che sale all'84,5\% all'ultimo checkpoint) rispetto al 33,3\% del caso;
(2)~la topologia del reticolo FCC apprende un accoppiamento tra canali
$4{,}6\times$ superiore rispetto alla topologia cubica, sebbene nessuna
delle due migliori la perdita di addestramento; (3)~l'interpolazione
lineare delle matrici ponte tra adattatori preserva la struttura
dell'autospettro (coseno $> 0{,}999$), consentendo un'analisi della
composizione degli adattatori in dimensione ridotta; (4)~il
pre-addestramento contrastivo della matrice ponte installa una
preferenza direzionale (rapporto co-planare/cross-planare $> 2.600\times$)
che persiste come memoria geometrica stabile dopo la rimozione della
supervisione---decadendo a $16\times$ in 6.000 passi non supervisionati
ma senza mai raggiungere la baseline nulla di $1{,}0$, mentre la
connettivit\`a complessiva della matrice ponte cresce indipendentemente
di $3{,}7\times$. La matrice ponte non migliora la perdita n\'e la
perplessi\`a rispetto al LoRA standard a parit\`a di rango. Il suo
contributo \`e diagnostico, non generativo: un riassunto a 36 parametri
della struttura di accoppiamento dell'adattatore, leggibile senza
inferenza n\'e valutazione.
\end{abstract}

% ===================================================================
% Section 1 — Introduzione
% ===================================================================

\section{Introduzione}
\label{sec:introduction}

Un adattatore LoRA addestrato~\cite{hu2022lora} \`e opaco.
Le sue due matrici di proiezione $A \in \mathbb{R}^{r \times d}$ e
$B \in \mathbb{R}^{d \times r}$ codificano quanto richiesto
dall'obiettivo di fine-tuning, ma non esiste un modo strutturato per
ispezionare \emph{cosa} abbiano appreso senza eseguire l'inferenza su
dati di valutazione. Due adattatori con architettura identica, addestrati
su task differenti, sono indistinguibili dalle sole statistiche dei pesi:
entrambi sono a basso rango, entrambi sono piccoli, entrambi producono
$\Delta W = BA$.
L'adattatore \`e una scatola nera dentro una scatola nera.

Questa opacit\`a ha conseguenze pratiche.
I metodi di fusione degli adattatori---aritmetica dei
task~\cite{ilharco2023editing},
TIES-Merging~\cite{yadav2023ties}, DARE~\cite{yu2024dare},
AdaMerging~\cite{yang2024adamerging}---operano su matrici dei pesi
complete o sui loro fattori SVD senza alcuna rappresentazione compatta
dell'identit\`a del task.
Il rilevamento del sovradattamento richiede curve di perdita sulla
validazione.
La selezione dell'adattatore da un hub o registro si basa sui metadati,
non sullo stato interno dell'adattatore.
In ciascun caso, l'adattatore stesso non fornisce alcun riassunto del
proprio contenuto.

Introduciamo una \emph{matrice ponte}: una matrice di accoppiamento
addestrabile $C \times C$ inserita tra $A$ e $B$ nella fattorizzazione
LoRA.
La proiezione discendente $A$ mappa l'input su un vettore nascosto di
rango $r$, che viene riorganizzato in $C$ canali di larghezza $r/C$;
la matrice ponte mescola i canali; il risultato viene riorganizzato e
passato alla proiezione ascendente $B$.
Con $C = 6$ (motivato dal numero di facce del dodecaedro rombico
nell'architettura RhombiLoRA~\cite{rhombilora_paper1, rhombilora_paper2}),
la matrice ponte aggiunge 36 parametri per livello dell'adattatore---un
sovraccarico dello $0{,}02\%$ rispetto ai parametri totali dell'adattatore.
Quando la matrice ponte \`e la matrice identit\`a $I_C$, l'architettura
si riduce esattamente al LoRA standard.

Attraverso sette esperimenti che coprono l'anatomia della matrice ponte,
il fingerprinting del task, la fusione degli adattatori e il
pre-addestramento contrastivo su Qwen2.5 (1,5B e 7B), stabiliamo
quattro risultati:

\paragraph{(a) Fingerprinting del task.}
Le matrici ponte codificano l'identit\`a del task.
Un SVM leave-one-out addestrato sulle 28 matrici ponte delle proiezioni
query (1.008 parametri totali) classifica il tipo di task con
un'accuratezza dell'84,5\%, rispetto al 33,3\% del caso ($p < 10^{-6}$,
Mann-Whitney $U$).
I task modulano la \emph{magnitudine} dell'accoppiamento della matrice
ponte---l'addestramento su istruzioni generali produce un accoppiamento
tra canali $2{,}2\times$ superiore rispetto alla generazione di codice
(connettivit\`a di Fiedler 0,040 contro 0,018)---preservando al contempo
la \emph{struttura} dell'accoppiamento: la similarit\`a del coseno
dell'autospettro supera 0,999 tra tutte le coppie di task.
L'impronta emerge in modo monotono dall'addestramento
(0\% al passo 0, 66\% al passo 500, 93\% al passo 1.500 per la coppia
binaria pi\`u difficile) senza alcun artefatto di inizializzazione,
poich\'e tutte le matrici ponte partono come $I_C$.

\paragraph{(b) Composizione regolare degli adattatori.}
L'interpolazione lineare delle matrici ponte tra adattatori specifici
per task preserva la struttura dell'autospettro (coseno $> 0{,}999$) a
tutti i rapporti di mescolamento $\alpha \in [0, 1]$, mentre le matrici
ponte con pesi casuali mantengono solo un coseno $\approx 0{,}5$.
La traiettoria del valore di Fiedler \`e lineare per coppie di task
strutturalmente simili (alpaca$\leftrightarrow$code, $R^2 = 0{,}948$)
ma non lineare per coppie dissimili (code$\leftrightarrow$math,
$R^2 = 0{,}782$), con interferenza distruttiva visibile come una
diminuzione della connettivit\`a a $\alpha \approx 0{,}2$---un fenomeno
invisibile alle diagnostiche di fusione a peso pieno.
Ci\`o suggerisce che l'interpolazione delle matrici ponte pu\`o servire
come sonda a bassa dimensionalit\`a della compatibilit\`a tra adattatori
prima di impegnarsi in una fusione a peso pieno costosa.

\paragraph{(c) Informativit\`a non uniforme.}
Le proiezioni query portano la struttura pi\`u discriminativa per il task
nella matrice ponte attraverso tutte e tre le fasi sperimentali:
accoppiamento pi\`u elevato (Fiedler 0,050 contro 0,035--0,039 per K/V/O),
distanza inter-task pi\`u elevata, sensibilit\`a alla fusione pi\`u
elevata.
La correlazione cross-fase di Pearson $r = 0{,}921$ tra la
discriminazione del task (Fase~1A) e la sensibilit\`a alla fusione
(Fase~2A) conferma che i moduli che codificano la maggiore identit\`a
del task sono anche i pi\`u vulnerabili all'interferenza inter-task
durante la composizione.
I livelli pi\`u profondi del trasformatore sono significativamente pi\`u
specifici per il task rispetto ai livelli iniziali (Spearman
$\rho = 0{,}701$, $p < 0{,}0001$), con la proiezione di uscita che
supera la proiezione query negli ultimi cinque livelli (Livello~24
$o_\text{proj}$: distanza 0,319, la singola posizione pi\`u
discriminativa nel modello).

\paragraph{(d) Memoria geometrica.}
Il pre-addestramento contrastivo della matrice ponte installa una
preferenza direzionale (rapporto co-planare/cross-planare
$> 2.600\times$ dopo 500 passi supervisionati) che persiste come
memoria geometrica stabile dopo la rimozione della supervisione:
il rapporto decade a $16\times$ in 6.000 passi non supervisionati
ma non raggiunge mai la baseline nulla di $1{,}0$ stabilita da due
esperimenti di controllo indipendenti. Durante lo stesso periodo non
supervisionato, la connettivit\`a complessiva della matrice ponte
cresce di $3{,}7\times$ indipendentemente---l'ottimizzatore scopre
autonomamente l'utilit\`a dell'accoppiamento tra canali, mantenendo
al contempo la struttura direzionale che gli \`e stata insegnata.

\paragraph{Il risultato nullo.}
La matrice ponte non migliora la perdita di addestramento n\'e la
perplessi\`a rispetto al LoRA standard a parit\`a di rango.
La motivazione originale---che la topologia di accoppiamento delle facce
del dodecaedro rombico avrebbe indotto un accoppiamento direzionale
specifico dei canali a partire da dati di addestramento
strutturati---ha prodotto un risultato nullo completo
(rapporto co-planare/cross-planare $= 1{,}002$, $p = 0{,}474$),
confermato da un controllo separato sul tasso di apprendimento della
matrice ponte ($10\times$ LR, Co/Cross $\approx 1{,}0$).
Le dimensioni di rango nel LoRA sono rotazionalmente simmetriche: per
qualsiasi $R$ ortogonale, $B R^\top R A = BA$.
Le etichette topologiche sui canali non sopravvivono a questa simmetria
senza supervisione esplicita.
La matrice ponte apprende \emph{che} mescolare i canali \`e utile
(la topologia FCC produce un accoppiamento tra canali $4{,}6\times$
superiore rispetto alla cubica), ma non \emph{quali} canali
preferire---a meno che non venga esplicitamente istruita.

\paragraph{Contributo.}
Non proponiamo un adattatore migliore.
Proponiamo una nuova \emph{diagnostica}: una matrice a 36 parametri che
riassume la struttura di accoppiamento scoperta dall'addestramento tra
$A$ e $B$, leggibile senza inferenza n\'e valutazione su dati di
riserva.
La matrice ponte fornisce un'impronta compatta del task (1.008 parametri
per il sottoinsieme ottimale delle proiezioni Q---pi\`u piccola di un
singolo vettore di stato nascosto in Qwen2.5-7B), uno spazio a bassa
dimensionalit\`a per l'analisi della composizione degli adattatori, e
una mappa per livello e per modulo di dove la specializzazione del task
si concentra nella rete.
A nostra conoscenza, nessun lavoro precedente utilizza una
rappresentazione intermedia strutturata all'interno di un adattatore
come diagnostica del comportamento dell'adattatore.

Il precedente architetturale pi\`u vicino \`e LoRA-XS~\cite{banaei2024loraxs},
che inserisce una matrice $r \times r$ tra $A$ e $B$ ma congela le
proiezioni esterne per minimizzare il numero di parametri.
Poich\'e $A$ e $B$ non si adattano al task in LoRA-XS, la matrice
intermedia non pu\`o servire come diagnostica di ci\`o che
l'apprendimento congiunto $A$-$B$ ha scoperto---essa \emph{\`e}
l'interezza dell'adattamento appreso.
RhombiLoRA addestra congiuntamente tutte e tre le componenti,
consentendo alla matrice ponte di catturare il pattern di accoppiamento
emergente piuttosto che portare l'intero carico dell'adattamento.
La Sezione~\ref{sec:related} discute questa distinzione e altri metodi
correlati in dettaglio.

% ===================================================================
% Section 2 — Architettura
% ===================================================================

\section{Architettura}
\label{sec:architecture}

\subsection{LoRA Standard}
\label{sec:standard-lora}

L'adattamento a basso rango~\cite{hu2022lora} congela una matrice dei
pesi pre-addestrata $W_0 \in \mathbb{R}^{d_\text{out} \times d_\text{in}}$
e apprende un aggiornamento a basso rango $\Delta W = BA$, dove
$B \in \mathbb{R}^{d_\text{out} \times r}$,
$A \in \mathbb{R}^{r \times d_\text{in}}$, e
$r \ll \min(d_\text{in}, d_\text{out})$.
Il passo in avanti adattato \`e
\begin{equation}
  h = W_0 x + \frac{\alpha}{r}\, B A x,
  \label{eq:lora}
\end{equation}
dove $\alpha$ \`e un fattore di scala fisso e
$x \in \mathbb{R}^{d_\text{in}}$.
$A$ \`e inizializzata tramite uniforme di Kaiming~\cite{he2015delving}
e $B$ \`e inizializzata a zero, cosicch\'e $\Delta W = 0$ all'inizio
dell'addestramento.
Il vettore nascosto di rango $r$ $z = Ax \in \mathbb{R}^r$ funge da
rappresentazione intermedia, ma LoRA non impone alcuna struttura su di
esso: le $r$ dimensioni sono una base arbitraria per un sottospazio a
basso rango.
Due dimensioni di $z$ interagiscono solo nella misura in cui $B$ le
mescola nella proiezione ascendente---non esiste un meccanismo esplicito
di accoppiamento all'interno dello spazio nascosto stesso.

\subsection{La Modifica con Matrice Ponte}
\label{sec:bridge}

Partizioniamo il vettore nascosto di rango $r$ in $C$ canali di
larghezza $r / C$ (richiedendo $C \mid r$) e inseriamo una \emph{matrice
ponte} addestrabile $M \in \mathbb{R}^{C \times C}$ che mescola i
canali prima della proiezione ascendente.
Concretamente, il passo in avanti adattato diventa:
\begin{equation}
  h = W_0 x + \frac{\alpha}{r}\, B\,
    \operatorname{vec}\!\bigl(M \cdot
    \operatorname{mat}_{C \times r/C}(Ax)\bigr),
  \label{eq:bridge}
\end{equation}
dove $\operatorname{mat}_{C \times r/C}$ riorganizza $\mathbb{R}^r
\to \mathbb{R}^{C \times (r/C)}$ e $\operatorname{vec}$ ne \`e
l'inversa.
Equivalentemente, in notazione indiciale: il vettore nascosto $z = Ax$
viene riorganizzato in $Z \in \mathbb{R}^{C \times (r/C)}$, la matrice
ponte agisce come $\tilde{Z}_{ik} = \sum_j M_{ij} Z_{jk}$, e
$\tilde{Z}$ viene appiattito di nuovo in $\mathbb{R}^r$ prima della
moltiplicazione per $B$.

Nei nostri esperimenti $r = 24$ e $C = 6$, cosicch\'e ogni canale porta
$r/C = 4$ fette di rango e $M$ \`e una matrice $6 \times 6$ con 36
parametri addestrabili.
In corrispondenza di un modulo target rappresentativo in Qwen2.5-7B
($d_\text{in} = d_\text{out} = 3.584$), le matrici di proiezione LoRA
standard contengono $2 \times 24 \times 3.584 = 172.032$ parametri; la
matrice ponte ne aggiunge 36, un sovraccarico dello $0{,}021\%$.

\paragraph{Recupero dell'identit\`a.}
Quando $M = I_C$, la sequenza riorganizzazione--moltiplicazione--appiattimento
\`e la mappa identit\`a su $\mathbb{R}^r$.
L'Equazione~\eqref{eq:bridge} si riduce esattamente
all'Equazione~\eqref{eq:lora}: l'architettura con matrice ponte \`e una
generalizzazione stretta del LoRA standard, con LoRA come caso speciale
$M = I_C$.
Inizializziamo $M = I_C$ all'inizio dell'addestramento cosicch\'e la
matrice ponte parta inerte e qualsiasi struttura di accoppiamento in $M$
a convergenza sia interamente appresa, non ereditata dall'inizializzazione.

\paragraph{Flusso del gradiente.}
Il gradiente della perdita rispetto a un elemento della matrice ponte
$M_{ij}$ \`e
\begin{equation}
  \frac{\partial \mathcal{L}}{\partial M_{ij}}
  = \sum_{k=1}^{r/C}
    \frac{\partial \mathcal{L}}{\partial \tilde{Z}_{ik}}
    \cdot Z_{jk},
  \label{eq:bridge-grad}
\end{equation}
dove $\tilde{Z}_{ik}$ riceve il suo gradiente attraverso $B$ e la
perdita a valle.
$M_{ij}$ accumula quindi un gradiente proporzionale alla
\emph{co-attivazione} tra il canale~$j$ (nell'uscita della proiezione
discendente) e il canale~$i$ (nel segnale di gradiente della proiezione
ascendente).
Se i dati di addestramento inducono correlazioni sistematiche tra
specifiche coppie di canali, le corrispondenti voci fuori diagonale di
$M$ crescono; se le attivazioni dei canali sono indipendenti, $M$
rimane vicino a $I_C$.
Questa selettivit\`a \`e il meccanismo che consente il fingerprinting
del task (\S\ref{sec:experiments}).

\paragraph{Simmetria rotazionale.}
Per qualsiasi $R \in \mathbb{R}^{r \times r}$ ortogonale, il prodotto
$BA$ \`e invariante sotto $B \mapsto BR^\top$, $A \mapsto RA$.
Le dimensioni di rango del LoRA standard sono dunque rotazionalmente
simmetriche---non esiste un intrinseco ``canale~1'' o ``canale~6''.
La matrice ponte non rompe questa simmetria: l'ottimizzatore pu\`o
permutare o ruotare i canali liberamente, e la matrice ponte
compenser\`a.
Ci\`o che la matrice ponte \emph{fornisce} \`e un riassunto leggibile
della struttura di accoppiamento scoperta dall'addestramento,
indipendentemente dalla rotazione scelta dall'ottimizzatore.
Due matrici ponte addestrate sullo stesso task non avranno voci
identiche, ma esibiranno le stesse propriet\`a spettrali (autovalori,
connettivit\`a di Fiedler) poich\'e queste sono invarianti per
rotazione.
Sfruttiamo questo nelle nostre analisi confrontando le matrici ponte
tramite i loro autospettri e la connettivit\`a grafo-teorica piuttosto
che la distanza voce per voce.

\subsection{Vincoli Topologici}
\label{sec:topology}

La matrice ponte $M$ \`e una matrice densa $C \times C$ per default.
Le sue voci fuori diagonale possono opzionalmente essere vincolate da
una maschera binaria $\mathcal{T} \in \{0, 1\}^{C \times C}$ derivata
da una topologia reticolare, imponendo $M_{ij} = 0$ ovunque
$\mathcal{T}_{ij} = 0$.
Durante l'addestramento, le voci mascherate non ricevono gradiente e
rimangono a zero.

Definiamo due maschere topologiche per $C = 6$, derivate dalla
connettivit\`a di reticoli tridimensionali:

\paragraph{Maschera FCC.}
I sei canali corrispondono alle sei coppie di facce antipodaliche del
dodecaedro rombico---equivalentemente, alle sei coppie di direzioni dei
primi vicini del reticolo cubico a facce centrate (FCC).
Due coppie di direzioni si \emph{accoppiano} quando condividono vertici
sul poliedro sottostante.
Nel dodecaedro rombico, le coppie co-planari (quelle giacenti nello
stesso piano coordinato, ad es.\ le direzioni $(1,1,0)$ e $(1,-1,0)$ nel
piano $xy$) condividono 4 vertici ottaedrici; le coppie cross-planari
ne condividono 2.
Tutte le coppie condividono almeno 2 vertici, quindi la maschera FCC \`e
il grafo completo su 6 nodi: $\mathcal{T}^\text{FCC}_{ij} = 1$ per
tutti gli $i, j$.
La maschera ammette 6 voci diagonali e $\binom{6}{2} = 15$ connessioni
fuori diagonale bidirezionali, per un totale di 36 voci non nulle.
In pratica, la matrice ponte non vincolata \`e la matrice ponte con
maschera FCC---la topologia consente l'accoppiamento completo.

\paragraph{Maschera cubica.}
Il reticolo cubico connette ogni nodo a 6 vicini lungo 3 direzioni
allineate agli assi, producendo 3 coppie di direzioni.
Per costruire una maschera cubica su 6 canali, manteniamo
l'accoppiamento solo all'interno di ogni coppia co-planare: i canali
$(0,1)$, $(2,3)$ e $(4,5)$ possono accoppiarsi, ma le voci
cross-planari sono azzerate.
Ci\`o produce una maschera a blocchi diagonali con tre blocchi
$2 \times 2$: $\mathcal{T}^\text{cubic}_{ij} = 1$ se
$\lfloor i/2 \rfloor = \lfloor j/2 \rfloor$, e 0 altrimenti.
La maschera cubica ha 6 voci diagonali e 3 connessioni fuori diagonale
bidirezionali (una per blocco), per un totale di 12 voci non nulle.

\paragraph{Relazione con MELoRA.}
Impostando la maschera a $\mathcal{T} = I_C$ (solo diagonale, nessun
accoppiamento fuori diagonale) si recupera la struttura a blocchi
diagonali di MELoRA~\cite{ren2024melora}, in cui ogni canale opera
indipendentemente.
La matrice ponte copre dunque un continuum: da canali completamente
indipendenti ($M = I_C$, MELoRA), attraverso l'accoppiamento parziale
(maschera cubica), fino all'accoppiamento completo (maschera FCC o
senza vincoli).
Tutti gli esperimenti in questo articolo utilizzano la matrice ponte
senza vincoli salvo diversa indicazione.

\subsection{Grandezze Diagnostiche}
\label{sec:diagnostics}

La matrice ponte \`e una matrice quadrata, reale e di piccole
dimensioni.
Ne estraiamo tre grandezze diagnostiche, ciascuna invariante rispetto
alla simmetria rotazionale delle dimensioni di rango
(\S\ref{sec:bridge}):

\paragraph{Connettivit\`a di Fiedler.}
Interpretiamo $|M_{ij}|$ per $i \neq j$ come il peso dell'arco tra i
canali $i$ e $j$ in un grafo pesato su $C$ nodi, e calcoliamo la
connettivit\`a algebrica (il secondo autovalore pi\`u piccolo del
Laplaciano pesato)~\cite{fiedler1973algebraic}.
Un valore di Fiedler pari a zero indica una matrice ponte disconnessa o
identit\`a; valori maggiori indicano un accoppiamento tra canali pi\`u
forte.
All'inizializzazione ($M = I_C$), Fiedler $= 0$.

\paragraph{Deviazione di Frobenius.}
$\|M - I_C\|_F$ misura quanto la matrice ponte si \`e allontanata
dalla sua inizializzazione identit\`a.
Si tratta di un riassunto scalare dell'accoppiamento totale appreso,
che combina sia la scalatura diagonale sia il mescolamento fuori
diagonale.

\paragraph{Autospettro.}
Il vettore degli autovalori $\lambda_1, \ldots, \lambda_C$ di $M$
(ordinati per modulo) fornisce un'impronta invariante per rotazione
della struttura di accoppiamento della matrice ponte.
Confrontiamo gli autospettri tra matrici ponte tramite similarit\`a del
coseno:
$\operatorname{cos}(\boldsymbol{\lambda}^{(a)},
\boldsymbol{\lambda}^{(b)})$.

\subsection{Descrizione della Figura}
\label{sec:arch-figure}

% [Figura 1: Diagramma dell'architettura — da renderizzare separatamente]
%
% La figura dovrebbe mostrare quanto segue, disposto da sinistra a
% destra come un diagramma di flusso dati:
%
% SINISTRA: Input x (una barra verticale etichettata d_in).
%
% BLOCCO 1: Proiezione discendente A (r × d_in), che produce
% z ∈ R^r (una barra verticale stretta etichettata r = 24).
%
% RIORGANIZZAZIONE: z viene riorganizzato in una griglia 6 × 4
% (6 righe = canali, 4 colonne = fette di rango per canale). Ogni
% riga è codificata per colore con una coppia di direzioni. Visivo:
% una piccola matrice/griglia con 6 righe colorate.
%
% BLOCCO 2: La matrice ponte M (6 × 6), mostrata come una piccola
% matrice quadrata sovrapposta alla griglia dei canali. Frecce o
% linee collegano ogni riga di canale alle righe di canale in uscita,
% con spessore della linea proporzionale a |M_{ij}|.
% L'inizializzazione identità è mostrata come versione trasparente
% (6 connessioni solo diagonali), con lo stato appreso che mostra
% connessioni fuori diagonale emergenti.
%
% RIORGANIZZAZIONE: La griglia 6 × 4 mescolata viene appiattita in R^r.
%
% BLOCCO 3: Proiezione ascendente B (d_out × r), che produce la
% correzione LoRA Δh ∈ R^{d_out}.
%
% DESTRA: La correzione viene sommata a W_0 x (mostrato come un
% percorso parallelo separato da x attraverso W_0, che si unisce a
% Δh tramite "+").
%
% RIQUADRO (in basso a destra): Due piccole mappe di calore 6 × 6
% affiancate:
% (a) M = I_6 (identità, all'inizializzazione — solo diagonale,
% tutto il fuori diagonale bianco), etichettata "Passo 0: LoRA standard";
% (b) una M appresa (che mostra struttura fuori diagonale con
% intensità variabile), etichettata "Passo 10K: accoppiamento appreso".
%
% ANNOTAZIONE: Un richiamo sul blocco della matrice ponte che recita
% "36 parametri (0,021% dell'adattatore)" con l'equazione
% M_{ij}: canale j → canale i.
%
% Didascalia: "Architettura RhombiLoRA. La matrice ponte M accoppia i
% C = 6 canali della rappresentazione nascosta di rango r tra le
% proiezioni discendente e ascendente A e B. All'inizializzazione
% M = I_6, recuperando il LoRA standard. L'addestramento scopre
% l'accoppiamento tra canali (riquadro), producendo una diagnostica
% compatta del comportamento dell'adattatore."

% ===================================================================
% Section 3 — Fondamenti Teorici
% ===================================================================

\section{Fondamenti Teorici}
\label{sec:theory}

L'utilit\`a diagnostica della matrice ponte si fonda su una semplice
osservazione sul flusso del gradiente. Il gradiente della perdita
rispetto all'elemento della matrice ponte $M_{ij}$ \`e proporzionale
alla co-attivazione tra il canale $j$ della proiezione discendente e il
canale $i$ del gradiente della proiezione ascendente
(Eq.~\ref{eq:bridge-grad}).
Le voci della matrice ponte accumulano quindi un riassunto continuo
della correlazione tra canali, modellata dall'obiettivo di
addestramento.

Tre propriet\`a del gradiente della matrice ponte la rendono
diagnostica:

\paragraph{Selettivit\`a.}
Se i dati di addestramento inducono correlazioni sistematiche tra
specifiche coppie di canali---poich\'e certe combinazioni di
caratteristiche co-occorrono tra gli esempi---le corrispondenti voci
fuori diagonale crescono. Se i canali sono indipendenti, $M$ rimane
vicino a $I_C$. La matrice ponte \`e modellata dalla struttura
statistica del task, non dall'inizializzazione n\'e dal bias
architetturale.

\paragraph{Compattezza.}
La matrice ponte occupa $C^2 = 36$ parametri per livello,
indipendentemente dalla dimensione del modello $d_\text{model}$.
Questo \`e pi\`u piccolo di un singolo vettore di stato nascosto nel
modello base. Il rapporto di compressione diagnostica \`e
$d_\text{model}^2 / C^2 \approx 3{,}6 \times 10^5$ per Qwen2.5-7B.

\paragraph{Diagnostiche invarianti.}
Sebbene le voci della matrice ponte dipendano dalla scelta
dell'ottimizzatore in merito alla rotazione nello spazio di rango,
l'autospettro di $M$ \`e invariante per rotazione. Due matrici ponte
addestrate sullo stesso task con semi casuali differenti differiranno
voce per voce ma convergeranno verso autospettri simili (similarit\`a
del coseno $> 0{,}999$ tra tutte le coppie di task nei nostri
esperimenti). L'analisi spettrale della matrice ponte \`e pertanto
riproducibile tra diverse esecuzioni di addestramento.

La previsione teorica---che la struttura del task fluisca nella
struttura della matrice ponte attraverso la co-attivazione del
gradiente---\`e confermata empiricamente nelle
Sezioni~\ref{sec:exp-fingerprint} e~\ref{sec:exp-merging}.

% ===================================================================
% Section 4 — Esperimenti
% ===================================================================

\section{Esperimenti}
\label{sec:experiments}

Valutiamo le matrici ponte di RhombiLoRA attraverso sette esperimenti,
procedendo dall'analisi strutturale di un singolo adattatore attraverso
la supervisione direzionale e la memoria geometrica fino al
fingerprinting multi-task e alla composizione degli adattatori a livello
di matrice ponte. Tutti gli esperimenti utilizzano la libreria
\texttt{rhombic} (v0.3.0) con pipeline identiche di estrazione delle
matrici ponte e analisi spettrale. Codice, matrici ponte e script di
riproduzione sono disponibili presso
\texttt{https://github.com/promptcrafted/rhombic}.

\paragraph{Configurazione comune.}
Salvo diversa indicazione, tutti gli adattatori utilizzano rango
$r = 24$ con $c = 6$ canali in topologia FCC, producendo matrici ponte
$6 \times 6$ (36 parametri addestrabili per adattatore).
L'addestramento utilizza QLoRA (quantizzazione NormalFloat a 4 bit) con
AdamW ($\beta_1 = 0{,}9$, $\beta_2 = 0{,}999$), tasso di apprendimento
$2 \times 10^{-4}$ con schedulazione a coseno, dimensione del batch 4,
accumulazione del gradiente 4, e precisione mista \texttt{bf16}. I
moduli target sono \texttt{q\_proj}, \texttt{k\_proj}, \texttt{v\_proj}
e \texttt{o\_proj} attraverso tutti i livelli del trasformatore.

% ---------------------------------------------------------------------------
\subsection{Anatomia della Matrice Ponte (1,5B)}
\label{sec:exp-anatomy-1b}

\paragraph{Configurazione.} Addestriamo cinque configurazioni di
adattatore su Qwen2.5-1.5B-Instruct (28 livelli, 112 adattatori
totali) utilizzando Alpaca-cleaned~\cite{alpaca} (52K esempi) per 2.000
passi su una NVIDIA RTX 6000 Ada (48\,GB). Le configurazioni sono:
(1)~LoRA standard $r\!=\!24$, (2)~RhombiLoRA $r\!=\!24$ con matrice
ponte identit\`a congelata, (3)~RhombiLoRA $r\!=\!24$ con matrice ponte
addestrabile (inizializzazione identit\`a), (4)~LoRA standard
$r\!=\!48$ (budget di parametri $2\times$), e (5)~RhombiLoRA
$r\!=\!24$ con topologia cubica ($c\!=\!3$, matrice ponte $3 \times 3$).
Le configurazioni 1--3 e~5 hanno conteggi di parametri addestrabili
corrispondenti ($\sim$3,27M); la configurazione~4 li raddoppia a
$\sim$6,54M. Le configurazioni di addestramento sono riassunte nella
Tabella~\ref{tab:exp1-configs}.

% Table: Exp 1 configurations
\begin{table}[t]
\centering
\caption{Configurazioni dell'Esperimento 1 (Qwen2.5-1.5B-Instruct, Alpaca-cleaned, 2K passi).}
\label{tab:exp1-configs}
\small
\begin{tabular}{@{}clcccc@{}}
\toprule
\# & Configurazione & Rango & Canali & Matrice Ponte & Parametri \\
\midrule
1 & LoRA Standard          & 24 & 6 & $I$ congelata & 3.268.608 \\
2 & RhombiLoRA (congelata)    & 24 & 6 & $I$ congelata & 3.268.608 \\
3 & RhombiLoRA (addestrabile) & 24 & 6 & Addestrabile  & 3.270.624 \\
4 & LoRA Standard          & 48 & 6 & $I$ congelata & 6.537.216 \\
5 & RhombiLoRA (cubica)     & 24 & 3 & Addestrabile  & 3.269.112 \\
\bottomrule
\end{tabular}
\end{table}

\paragraph{Risultati.}
\textbf{Equivalenza dell'identit\`a.} Le configurazioni 1 e~2 producono
traiettorie di perdita con $\Delta < 0{,}0002$ a ogni checkpoint,
confermando che RhombiLoRA con matrice ponte identit\`a congelata
recupera esattamente il LoRA standard.

\textbf{Saturazione del rango.} La configurazione~4 (rango 48,
$2\times$ parametri) converge alla stessa perdita finale del rango~24
(0,376 contro 0,376 al passo 2000). Il collo di bottiglia del segnale
di addestramento a 2K passi non \`e la capacit\`a di rango; la matrice
ponte fornisce un tipo qualitativamente differente di capacit\`a
(mescolamento tra canali) piuttosto che una maggiore quantit\`a dello
stesso tipo.

\textbf{Accoppiamento strutturato.} La matrice ponte della
configurazione~3 sviluppa una connettivit\`a algebrica (valore di
Fiedler) in crescita monotona durante l'addestramento: da
$\lambda_2 = 0{,}000$ all'inizializzazione a $\lambda_2 = 0{,}037$ al
passo 2000. Il comportamento primario appreso \`e
l'auto-amplificazione diagonale ($\sim$$+$0,009 per canale);
l'accoppiamento fuori diagonale tra canali \`e di un ordine di grandezza
inferiore ($\sim$0,001) ma costantemente presente e in crescita.

\textbf{La topologia conta.} Il risultato critico: la matrice ponte FCC
a 6 canali sviluppa un valore di Fiedler \textbf{4,6$\times$} superiore
e una deviazione dall'identit\`a \textbf{5,3$\times$} superiore rispetto
alla matrice ponte cubica a 3 canali
(Tabella~\ref{tab:topology-effect}), con budget di parametri
corrispondenti e perdita finale identica. Entrambe le topologie
sviluppano la stessa auto-amplificazione diagonale ($+$0,009), ma la
topologia FCC scopre e sfrutta pi\`u percorsi di accoppiamento tra
canali. L'ottimizzatore trova che 12 archi (FCC) offrono una struttura
di mescolamento pi\`u utile di 6 archi (cubica)---confermando a livello
di rete neurale il vantaggio di connettivit\`a algebrica stabilito dai
benchmark grafo-teorici~\cite{rhombilora_paper1}.

\begin{table}[t]
\centering
\caption{Effetto della topologia sulla struttura della matrice ponte
(Esp.\ 1, passo 2000). Budget di parametri identico, perdita identica.
La topologia FCC sviluppa una connettivit\`a algebrica $4{,}6\times$
superiore.}
\label{tab:topology-effect}
\small
\begin{tabular}{@{}lccc@{}}
\toprule
Metrica & FCC (6 canali) & Cubica (3 canali) & Rapporto \\
\midrule
Crescita diagonale         & $+$0,0087 & $+$0,0088 & 1,0$\times$ \\
Media $|$fuori diag.$|$ & 0,00100   & 0,00060   & 1,7$\times$ \\
Fiedler $\lambda_2$     & 0,0373    & 0,0082    & 4,6$\times$ \\
Deviazione da $I$      & 0,0507    & 0,0095    & 5,3$\times$ \\
\bottomrule
\end{tabular}
\end{table}

% ---------------------------------------------------------------------------
\subsection{Anatomia della Matrice Ponte a Scala (7B)}
\label{sec:exp-anatomy-7b}

\paragraph{Configurazione.} Addestriamo RhombiLoRA ($r\!=\!24$, FCC,
matrice ponte addestrabile) su Qwen2.5-7B-Instruct (28 livelli, 4
moduli per livello, 112 adattatori) utilizzando Alpaca-cleaned per
10.000 passi su una RTX 6000 Ada locale. Le matrici ponte sono estratte
a 11 checkpoint (passi 0, 1K, 2K, \ldots, 10K), producendo 1.344
istantanee della matrice ponte. Analizziamo le 112 matrici ponte finali
per la struttura posizionale e la traiettoria completa per le dinamiche
di addestramento.

\paragraph{Struttura per modulo.}
L'accoppiamento della matrice ponte varia significativamente per tipo di
modulo di attenzione. L'ANOVA a una via sui valori finali di Fiedler
produce $F = 6{,}022$, $p = 0{,}0008$
(Tabella~\ref{tab:module-anatomy}). Le proiezioni Q sviluppano il
maggior accoppiamento (media $\lambda_2 = 0{,}050$), il 40\% in pi\`u
rispetto alle proiezioni V ($\lambda_2 = 0{,}036$). L'ordinamento
Q $>$ K $>$ V non segue alcun pattern banalmente prevedibile.
In modo notevole, le strutture della matrice ponte Q-V sono pi\`u simili
di quelle Q-K ($t$ accoppiato $= 2{,}13$, $p = 0{,}042$), contrariamente
all'aspettativa ingenua che il prodotto scalare query-chiave legherebbe
Q e K pi\`u strettamente. La matrice ponte traccia il \emph{flusso di
informazione} (Q interroga ci\`o che V fornisce) piuttosto che il
\emph{calcolo di similarit\`a} (Q$\cdot$K).

\begin{table}[t]
\centering
\caption{Statistiche della matrice ponte per tipo di modulo (Esp.\ 2,
Qwen2.5-7B, passo 10K, $n = 28$ per modulo). ANOVA: $F = 6{,}022$,
$p = 0{,}0008$.}
\label{tab:module-anatomy}
\small
\begin{tabular}{@{}lcccc@{}}
\toprule
Modulo & Fiedler ($\lambda_2$) & Deviazione & Disp.\ Autov. & Co/Cross \\
\midrule
q\_proj & $0{,}050 \pm 0{,}014$ & $0{,}273 \pm 0{,}070$ & $0{,}124 \pm 0{,}025$ & $1{,}050 \pm 0{,}593$ \\
k\_proj & $0{,}039 \pm 0{,}014$ & $0{,}155 \pm 0{,}034$ & $0{,}108 \pm 0{,}029$ & $1{,}074 \pm 0{,}528$ \\
o\_proj & $0{,}035 \pm 0{,}016$ & $0{,}225 \pm 0{,}099$ & $0{,}101 \pm 0{,}044$ & $1{,}096 \pm 0{,}411$ \\
v\_proj & $0{,}036 \pm 0{,}016$ & $0{,}142 \pm 0{,}064$ & $0{,}094 \pm 0{,}033$ & $1{,}061 \pm 0{,}448$ \\
\bottomrule
\end{tabular}
\end{table}

\paragraph{Invarianza rispetto alla profondit\`a.}
Contrariamente al tipo di modulo, la profondit\`a del livello non ha
effetto significativo sui valori di Fiedler (ANOVA $F = 0{,}799$,
$p = 0{,}74$; Pearson $r = 0{,}03$, $p = 0{,}87$).
La matrice ponte sviluppa una connettivit\`a algebrica comparabile a
ogni profondit\`a. Tuttavia, la \emph{deviazione} dall'identit\`a
aumenta con la profondit\`a (ANOVA $F = 2{,}43$, $p = 0{,}001$): i
livelli pi\`u profondi divergono di pi\`u dall'inizializzazione
identit\`a, suggerendo che, sebbene tutti i livelli sviluppino una forza
di accoppiamento simile, i livelli pi\`u profondi sviluppano pattern di
accoppiamento pi\`u complessi (asimmetrici o non uniformi).

\paragraph{Dinamiche di addestramento.}
Il valore medio di Fiedler cresce in modo monotono lungo l'intera
traiettoria: $0{,}000 \to 0{,}018$ (1K) $\to 0{,}040$ (10K). La
correlazione anticipata--tardiva \`e debole ($r = 0{,}099$,
$p = 0{,}30$): gli adattatori che sviluppano il maggior accoppiamento
al passo 1K non mantengono necessariamente quella posizione al passo
10K. La matrice ponte subisce una riorganizzazione strutturale durante
l'addestramento piuttosto che un'amplificazione uniforme della struttura
iniziale. Il coefficiente di variazione tra le matrici ponte finali \`e
$\text{CV} = 0{,}408$, con un intervallo di $9\times$ tra l'adattatore
pi\`u accoppiato (Livello~2 \texttt{q\_proj}, $\lambda_2 = 0{,}099$) e
il meno accoppiato (Livello~1 \texttt{v\_proj},
$\lambda_2 = 0{,}011$).

\paragraph{Implicazione.}
La matrice ponte \`e \emph{architetturalmente sensibile} alla sua
posizione all'interno della testa di attenzione: il tipo di modulo
spiega una varianza significativa nella forza di accoppiamento.
Ci\`o motiva l'ipotesi che le matrici ponte possano anche essere
\emph{sensibili al task}---se una matrice a 36 parametri pu\`o
distinguere \texttt{q\_proj} da \texttt{v\_proj}, potrebbe distinguere
``codice'' da ``matematica''.
Verifichiamo ci\`o direttamente in \S\ref{sec:exp-fingerprint}.

% ---------------------------------------------------------------------------
\subsection{Direzione vs.\ Connettivit\`a}
\label{sec:exp-direction}

\paragraph{Motivazione.}
La topologia della matrice ponte FCC deriva dal dodecaedro rombico, le
cui sei coppie di canali si decompongono in tre coppie co-planari
(che condividono 4 vertici ottaedrici) e tre coppie cross-planari
(che ne condividono 2). Abbiamo ipotizzato che dati di addestramento con
struttura esplicita di accoppiamento prospettico avrebbero indotto un
accoppiamento co-planare misurabilmente pi\`u forte di quello
cross-planare nella matrice ponte.

\paragraph{Configurazione.}
Addestriamo RhombiLoRA ($r\!=\!24$, FCC) su Qwen2.5-7B-Instruct
utilizzando un dataset geometrico sintetico (23K esempi con il 95\%
di accoppiamento prospettico co-planare) per 3.000 passi su una RunPod
RTX 4090. Confrontiamo con la baseline Alpaca (Esp.~2, dati isotropi,
10K passi) utilizzando quattro diagnostiche spettrali: rapporto di
accoppiamento co-planare/cross-planare, test di permutazione,
allineamento della matrice di Gram e ritenzione della perturbazione.

\paragraph{Risultato nullo sulla direzione.}
Il dataset geometrico non produce \emph{alcun} segnale direzionale. Il
rapporto co/cross \`e 1,002 al passo 3K (effettivamente identit\`a),
rispetto al picco transitorio di Alpaca di 1,091 al passo 3K che decade
a 1,019 al passo 10K. Nessuno dei due risultati \`e statisticamente
significativo sotto test di permutazione (geometrico $p = 0{,}474$;
Alpaca $p = 0{,}332$). Il confronto completo \`e presentato nella
Tabella~\ref{tab:direction-null}.

\begin{table}[t]
\centering
\caption{Confronto del segnale direzionale (112 matrici ponte FCC
ciascuno). I dati geometrici sono stati progettati con il 95\% di
accoppiamento prospettico co-planare. Nessun dataset produce una
preferenza direzionale significativa.}
\label{tab:direction-null}
\small
\begin{tabular}{@{}lccl@{}}
\toprule
Metrica & Alpaca (10K) & Geometrico (3K) & Migliore \\
\midrule
Rapporto Co/Cross         & 1,019 & 1,002 & --- \\
Permutazione $p$        & 0,332 & 0,474 & --- \\
Distanza di Gram          & 0,921 & 0,921 & Parit\`a \\
Dispersione autovalori      & 0,107 & 0,062 & Alpaca \\
Ritenzione perturbazione & 68,8\% & 69,7\% & Parit\`a \\
Fiedler $\lambda_2$    & 0,040 & 0,030 & Alpaca \\
\bottomrule
\end{tabular}
\end{table}

\paragraph{Risultato positivo sulla connettivit\`a.}
Mentre il risultato nullo sulla direzione \`e definitivo, la struttura
della matrice ponte \`e \emph{distribuita} indipendentemente dal tipo di
dati: azzerando tutte le voci co-planari della matrice ponte si
mantiene $\sim$69\% del valore di Fiedler in entrambe le condizioni
(68,8\% contro 69,7\%). Questo \`e un invariante architetturale, non una
propriet\`a dipendente dai dati. La matrice ponte distribuisce il suo
accoppiamento appreso sia sulle coppie co-planari sia su quelle
cross-planari indipendentemente dalla struttura dell'input.

\paragraph{Causa principale.}
Il bias co-planare a livello di prompt (95\% di accoppiamento
prospettico) non si traduce in una preferenza di co-attivazione a
livello di canale perch\'e \emph{l'assegnazione dei canali \`e
arbitraria}. La matrice ponte non ha alcun meccanismo per associare
prospettive semantiche (ad es.\ ``analitica + empirica'') con indici di
canale specifici (ad es.\ canali 0--1 contro 2--3). Ventotto livelli di
trasformatore tra la struttura del prompt e i pesi della matrice ponte
cancellano il segnale a livello di prompt. Indurre una preferenza
direzionale richiede la supervisione esplicita della struttura di
accoppiamento della matrice ponte---verifichiamo ci\`o direttamente in
\S\ref{sec:exp-contrastive}.

% ---------------------------------------------------------------------------
\subsection{Pre-addestramento Contrastivo della Matrice Ponte}
\label{sec:exp-contrastive}

\paragraph{Motivazione.}
Il risultato nullo in \S\ref{sec:exp-direction} ha stabilito che
l'assegnazione dei canali \`e arbitraria: la matrice ponte non ha alcun
meccanismo per associare canali specifici a direzioni specifiche del
reticolo a partire dai soli dati. Verifichiamo se una
\emph{supervisione esplicita}---una perdita contrastiva che premia
l'accoppiamento co-planare e penalizza quello cross-planare---possa
installare una preferenza direzionale, e se tale preferenza persista
dopo la rimozione della supervisione.

\paragraph{Configurazione.}
Addestriamo RhombiLoRA ($r\!=\!24$, FCC, matrice ponte addestrabile)
su Qwen2.5-7B-Instruct utilizzando Alpaca-cleaned in due fasi su una
RunPod RTX 3090:

\begin{enumerate}
\item \textbf{Riscaldamento contrastivo} (500 passi): La perdita di
addestramento
$\mathcal{L} = \mathcal{L}_\text{LM} + w_c \mathcal{L}_\text{contrastive}$
combina l'obiettivo standard di modellazione linguistica con un termine
contrastivo che massimizza l'accoppiamento tra coppie di canali
co-planari (quelle che condividono 4 vertici ottaedrici) rispetto alle
coppie cross-planari (che ne condividono 2). Il peso contrastivo cresce
linearmente da 0 a $w_c = 0{,}01$ in 500 passi.

\item \textbf{Continuazione non supervisionata} (passi 501--10.000+):
La perdita contrastiva viene interamente rimossa ($w_c = 0$).
L'addestramento prosegue con il solo obiettivo di modellazione
linguistica. La matrice ponte riceve gradienti unicamente attraverso il
percorso standard
$\partial\mathcal{L}_\text{LM}/\partial M_{ij}$.
\end{enumerate}

\noindent
Le matrici ponte sono salvate ogni 100 passi. Il rapporto di
accoppiamento co-planare/cross-planare (Co/Cross) fornisce la
diagnostica principale: a quale velocit\`a il segnale direzionale
supervisionato decade sotto addestramento non supervisionato, e
raggiunge la baseline nulla di 1,0 stabilita dagli Esperimenti~2.5
e~2.7?

\paragraph{Riscaldamento contrastivo.}
Al passo 500, la matrice ponte esibisce una massiccia preferenza
direzionale: Co/Cross $= 2.660\times$
(Tabella~\ref{tab:geometric-memory}). La perdita contrastiva installa
con successo la semantica geometrica che l'addestramento guidato dai
dati (\S\ref{sec:exp-direction}) non era riuscito a produrre: le coppie
co-planari portano ordini di grandezza pi\`u peso di accoppiamento
rispetto alle coppie cross-planari.

\paragraph{Memoria geometrica.}
Dopo la rimozione della supervisione al passo 500, il segnale
direzionale decade rapidamente ma \emph{non raggiunge mai la baseline
nulla} (Tabella~\ref{tab:geometric-memory}). Al passo 5.000 (4.500
passi non supervisionati), il rapporto Co/Cross \`e decaduto da
$2.660\times$ a $18{,}6\times$---una riduzione del 99,3\% in
magnitudine, ma comunque $18{,}6\times$ al di sopra della baseline nulla
di $1{,}0$ stabilita da due esperimenti indipendenti
(\S\ref{sec:exp-direction}, \S\ref{sec:exp-bridge-lr}). Al passo 6.500,
il rapporto si stabilizza intorno a $16\times$, con decrementi
successivi su 500 passi che si riducono sotto 0,5.

\begin{table}[t]
\centering
\caption{Memoria geometrica: rapporto di accoppiamento
co-planare/cross-planare (Co/Cross), valore di Fiedler e perdita di
validazione dopo 500 passi di riscaldamento contrastivo seguiti da
addestramento non supervisionato. Il segnale direzionale decade verso la
baseline nulla ($\approx$1,0) ma non la raggiunge mai.
L'addestramento era in corso al momento della stesura; la tendenza
asintotica \`e gi\`a chiara.}
\label{tab:geometric-memory}
\small
\begin{tabular}{@{}rlrrr@{}}
\toprule
Passo & Fase & Co/Cross & Fiedler $\lambda_2$ & Perd.\ Val. \\
\midrule
0     & Iniz.            & 1,000    & 0,000 & --- \\
500   & Fine risc.      & 2.660  & ---   & --- \\
3.500 & Non superv.  & 56,0   & 0,007 & --- \\
5.000 & Non superv.  & 18,6   & 0,022 & 0,301 \\
5.500 & Non superv.  & 17,3   & 0,024 & 0,297 \\
6.500 & Non superv.  & 16,1   & 0,026 & 0,295 \\
\midrule
\multicolumn{2}{@{}l}{\textit{Baseline nulle}} \\
\quad Esp.\ 2.5 & Dati geometrici    & 1,002  & 0,030 & --- \\
\quad Esp.\ 2.7 & LR ponte elevato    & $\approx$1,0 & --- & --- \\
\bottomrule
\end{tabular}
\end{table}

\paragraph{La connettivit\`a emerge indipendentemente.}
Mentre il segnale direzionale decade, la connettivit\`a
\emph{complessiva} della matrice ponte cresce in modo monotono. Il
valore di Fiedler aumenta da $0{,}007$ al passo 3.500 a $0{,}026$ al
passo 6.500---un incremento di $3{,}7\times$ guidato interamente da
gradienti non supervisionati. La deviazione dall'identit\`a aumenta
analogamente ($0{,}151 \to 0{,}193$), indicando che la matrice ponte
diventa \emph{pi\`u} espressiva, non meno, durante la fase non
supervisionata. La perdita di validazione migliora costantemente
($0{,}307 \to 0{,}295$). La matrice ponte si sta riorganizzando---cedendo
l'eccesso di specificit\`a direzionale in favore di un accoppiamento
utile al task---mantenendo al contempo una memoria geometrica persistente
della struttura direzionale mostratale durante il riscaldamento.

\paragraph{Risultato rivisto.}
Il risultato originale L-001 (``la matrice ponte non pu\`o apprendere la
preferenza direzionale dai soli dati,'' \S\ref{sec:exp-direction}) deve
essere rivisto: \emph{la matrice ponte non pu\`o scoprire la direzione
dai soli dati, ma una volta mostrata la direzione attraverso supervisione
esplicita, mantiene una memoria geometrica stabile che migliaia di passi
di addestramento non supervisionato non riescono a cancellare.} La
direzione richiede insegnamento; una volta insegnata, persiste.

% ---------------------------------------------------------------------------
\subsection{Controllo del Tasso di Apprendimento della Matrice Ponte}
\label{sec:exp-bridge-lr}

\paragraph{Configurazione.}
Per escludere le dinamiche di ottimizzazione come spiegazione
alternativa del risultato nullo direzionale
(\S\ref{sec:exp-direction}), addestriamo RhombiLoRA ($r\!=\!24$, FCC,
matrice ponte addestrabile) su Qwen2.5-7B-Instruct utilizzando
Alpaca-cleaned per 3.000 passi su una RunPod RTX 3090, con un tasso di
apprendimento separato per i parametri della matrice ponte fissato a
$10\times$ il tasso base dell'adattatore ($2 \times 10^{-3}$ contro
$2 \times 10^{-4}$). Non viene applicata alcuna perdita contrastiva. Se
il risultato nullo direzionale nell'Esperimento~2.5 fosse dovuto a un
tasso di apprendimento insufficiente della matrice ponte---la matrice
ponte semplicemente non aggiornandosi abbastanza velocemente rispetto
alle matrici di proiezione---allora amplificare il gradiente della
matrice ponte dovrebbe produrre un segnale direzionale.

\paragraph{Risultato nullo.}
Il rapporto co-planare/cross-planare rimane a $\approx$1,0 per tutto
l'addestramento---indistinguibile dall'inizializzazione identit\`a e
dal risultato nullo dell'Esperimento~2.5
(Tabella~\ref{tab:geometric-memory}, righe inferiori).
Aumentare il tasso di apprendimento della matrice ponte di $10\times$
non induce preferenza direzionale. Combinato con il risultato
dell'Esperimento~2.5 (LR standard, dati geometrici) e il risultato
dell'Esperimento~2.6 (LR standard, supervisione contrastiva), ci\`o
stabilisce che la preferenza direzionale richiede \emph{supervisione
contrastiva esplicita}, non un'ottimizzazione pi\`u rapida n\'e dati
strutturati.

% ---------------------------------------------------------------------------
\subsection{Fingerprinting del Task}
\label{sec:exp-fingerprint}

\paragraph{Configurazione.}
Addestriamo tre adattatori specifici per task su Qwen2.5-7B-Instruct:
(A)~Alpaca-cleaned (seguire istruzioni generali, 52K esempi, 10K passi,
RTX 6000 Ada locale), (B)~CodeAlpaca-20k (generazione di codice, 20K
esempi, 2K passi, RTX 6000 Ada locale), e (C)~GSM8K (ragionamento
matematico, 7,5K esempi, 2K passi, RunPod RTX 4090). Tutti utilizzano
una configurazione RhombiLoRA identica ($r\!=\!24$, FCC, matrice ponte
addestrabile). Le matrici ponte sono estratte all'ultimo checkpoint,
producendo 336 matrici ponte totali (112 per task, 28 livelli
$\times$ 4 moduli). Inoltre, le matrici ponte sono estratte a
checkpoint intermedi (passi 0, 500, 1000, 1500) per gli adattatori di
codice e matematica per tracciare l'emergenza dell'impronta.

\paragraph{Distanza inter-task vs.\ intra-task.}
Le distanze delle matrici ponte (norma di Frobenius media su 112
posizioni) tra i task sono significativamente maggiori delle distanze
intra-task. Mann-Whitney $U = 3.977.447$, $p < 10^{-6}$
(Tabella~\ref{tab:fingerprint-distances}).
Le tre coppie di task non sono equidistanti: alpaca--codice (0,189) e
alpaca--matematica (0,171) sono ben separate, mentre codice--matematica
(0,066) \`e sostanzialmente pi\`u vicina. Questa asimmetria \`e
interpretabile: codice e matematica condividono pattern di ragionamento
strutturato e formale che distinguono entrambi dal dominio diversificato
delle istruzioni generali.

\begin{table}[t]
\centering
\caption{Statistiche della distanza delle matrici ponte (Fase~1A). Le
distanze inter-task superano significativamente le distanze intra-task
(Mann-Whitney $p < 10^{-6}$).}
\label{tab:fingerprint-distances}
\small
\begin{tabular}{@{}lrrrr@{}}
\toprule
Coppia & Media & Dev.\ St. & Min & Max \\
\midrule
\multicolumn{5}{@{}l}{\textit{Inter-task}} \\
\quad alpaca $\leftrightarrow$ codice & 0,189 & 0,087 & 0,025 & 0,443 \\
\quad alpaca $\leftrightarrow$ matem. & 0,171 & 0,078 & 0,039 & 0,441 \\
\quad codice $\leftrightarrow$ matem.   & 0,066 & 0,029 & 0,024 & 0,170 \\
\midrule
\multicolumn{5}{@{}l}{\textit{Intra-task}} \\
\quad alpaca & 0,177 & 0,061 & --- & --- \\
\quad codice   & 0,047 & 0,015 & --- & --- \\
\quad matem.   & 0,085 & 0,029 & --- & --- \\
\bottomrule
\end{tabular}
\end{table}

\paragraph{Accuratezza della classificazione.}
La classificazione SVM leave-one-out sui 336 vettori delle matrici ponte
raggiunge il 72,9\% di accuratezza a tre vie (caso: 33,3\%). La
prestazione non \`e uniforme: alpaca (93,8\%) e codice (94,6\%) sono ben
classificati, mentre la matematica (50,0--66,1\%) \`e frequentemente
confusa con il codice, coerentemente con la loro bassa distanza
inter-task.

\paragraph{Selezione delle caratteristiche: dominanza della proiezione Q.}
Restringendo il classificatore alle sole matrici ponte della proiezione
Q (28 per task, 1.008 parametri) l'accuratezza \emph{migliora}
all'82,1\% (Tabella~\ref{tab:subset-class}).
Aggiungendo le matrici ponte della proiezione O (56 totali, 2.016
parametri) si ottiene la migliore accuratezza complessiva dell'83,3\%,
principalmente migliorando la classificazione della matematica dal
50,0\% al 66,1\%. Le matrici ponte delle proiezioni K e V non sono
informative: la loro inclusione degrada il confine decisionale del SVM.
Il risultato principale \`e che \textbf{1.008 parametri delle matrici
ponte (28 ponti Q-proj) superano l'intero insieme di 4.032 parametri}
per l'identificazione del task. Le matrici ponte della proiezione query
sono la componente pi\`u discriminativa per il task nella testa di
attenzione, coerentemente con la loro forza di accoppiamento pi\`u
elevata (\S\ref{sec:exp-anatomy-7b}).

\begin{table}[t]
\centering
\caption{Accuratezza della classificazione per sottoinsieme (LOO SVM, 3
vie). La rimozione delle matrici ponte K/V non informative migliora
l'accuratezza fino a 10,4\,pp. Le sole matrici ponte Q-proj
classificano all'82,1\% con 1.008 parametri.}
\label{tab:subset-class}
\small
\begin{tabular}{@{}lrrr@{}}
\toprule
Sottoinsieme & Posizioni & Parametri & Accuratezza LOO \\
\midrule
Completo (Q+K+V+O) & 112 & 4.032 & 72,9\% \\
Solo Q-proj     &  28 & 1.008 & 82,1\% \\
Solo O-proj     &  28 & 1.008 & 76,2\% \\
Q+O combinati    &  56 & 2.016 & \textbf{83,3\%} \\
\bottomrule
\end{tabular}
\end{table}

\paragraph{Gradiente di profondit\`a per livello.}
La discriminazione del task aumenta con la profondit\`a del livello.
Spearman $\rho = 0{,}701$ ($p < 0{,}0001$) tra l'indice del livello e
la distanza media inter-task di Frobenius. I primi 5 livelli pi\`u
discriminativi si trovano tutti nell'ultimo quarto della rete (livelli
20--24), con un intervallo dinamico di $10{,}3\times$ tra la posizione
pi\`u discriminativa (Livello~24 \texttt{o\_proj}, distanza $= 0{,}319$)
e la meno discriminativa (Livello~1 \texttt{v\_proj},
distanza $= 0{,}031$).
Le proiezioni di uscita dei livelli tardivi superano le proiezioni query
nelle singole posizioni---la proiezione di uscita si specializza per
task nei livelli in cui il modello sta formattando il suo output finale.

\paragraph{Magnitudine, non struttura.}
I task modificano la \emph{magnitudine} dell'accoppiamento della matrice
ponte, non la \emph{struttura} dell'accoppiamento. La similarit\`a del
coseno dell'autospettro tra coppie di task \`e $\geq 0{,}9997$ in tutte
le 112 posizioni. La forma spettrale della matrice ponte \`e invariante
tra i task; solo la scala differisce. Gli adattatori alpaca sviluppano
il Fiedler medio pi\`u elevato ($0{,}040$), seguiti dalla matematica
($0{,}028$) e dal codice ($0{,}018$). Il seguito di istruzioni
diversificate richiede pi\`u mescolamento tra canali rispetto ai domini
formali e focalizzati.

\paragraph{Emergenza dell'impronta.}
Le impronte del task emergono interamente dall'addestramento. Al
passo~0, tutte le matrici ponte sono identiche (inizializzazione
identit\`a) e tutte le distanze inter-task sono esattamente 0,000. La
divergenza \`e monotona: al passo~1000 il classificatore binario
codice--matematica raggiunge l'87,5\%, salendo al 92,9\% al passo~1500
e ancora in crescita. Il classificatore a tre vie sulle proiezioni Q
raggiunge l'84,5\% all'ultimo checkpoint. Questa traiettoria esclude
qualsiasi artefatto di inizializzazione e dimostra che la struttura
della matrice ponte \`e una registrazione fedele della specializzazione
indotta dall'addestramento.

% ---------------------------------------------------------------------------
\subsection{Fusione degli Adattatori a Livello di Matrice Ponte}
\label{sec:exp-merging}

\paragraph{Configurazione.}
Interpoliamo le matrici ponte finali tra tutte e tre le coppie di task a
11 coefficienti di mescolamento uniformemente spaziati ($\alpha \in
\{0{,}0,\; 0{,}1,\; \ldots,\; 1{,}0\}$), dove
$B_\alpha = (1 - \alpha) B_A + \alpha B_B$. Per ogni $\alpha$,
calcoliamo il valore di Fiedler, la deviazione dall'identit\`a, la
distanza di Frobenius da ciascuna sorgente e la similarit\`a del coseno
dell'autospettro rispetto a ciascuna sorgente su tutte le 112 posizioni
degli adattatori. Confrontiamo con una baseline casuale: fusione di ogni
matrice ponte del task con 20 matrici $6 \times 6$ casuali indipendenti
a $\alpha = 0{,}5$. Non \`e necessaria alcuna GPU---l'intero
esperimento opera su matrici a 36 parametri tramite algebra lineare.

\paragraph{Traiettorie di interpolazione.}
Le tre coppie di task esibiscono comportamenti di fusione
qualitativamente differenti (Tabella~\ref{tab:merge-linearity}).
L'interpolazione alpaca--codice \`e approssimativamente lineare
($R^2 = 0{,}948$): il Fiedler decresce regolarmente dall'estremo alpaca
(0,040) all'estremo codice (0,018) con residuo massimo 0,0035. Al
contrario, alpaca--matematica ($R^2 = 0{,}735$) e codice--matematica
($R^2 = 0{,}782$) esibiscono traiettorie non lineari. La coppia
codice--matematica mostra una \emph{diminuzione del Fiedler} a
$\alpha \approx 0{,}2$: la connettivit\`a scende da 0,018 a 0,015
prima di recuperare a 0,028, indicando interferenza distruttiva a
rapporti di mescolamento intermedi tra adattatori strutturalmente
dissimili.

\begin{table}[t]
\centering
\caption{Linearit\`a dell'interpolazione e preservazione della struttura
della matrice ponte. La non linearit\`a correla con la dissimilarit\`a
strutturale tra i task. Il coseno dell'autospettro $> 0{,}999$ a tutti
gli $\alpha$ per le fusioni tra task contro $\sim$0,5 per le baseline
casuali.}
\label{tab:merge-linearity}
\small
\begin{tabular}{@{}lccc@{}}
\toprule
Coppia & $R^2$ (Fiedler) & Res.\ Max. & $\cos$ autov.\ @ $\alpha\!=\!0{,}5$ \\
\midrule
alpaca $\leftrightarrow$ codice & 0,948 & 0,0035 & 0,9999 \\
alpaca $\leftrightarrow$ matem. & 0,735 & 0,0044 & 0,9999 \\
codice $\leftrightarrow$ matem.   & 0,782 & 0,0035 & 1,0000 \\
\midrule
\textit{Baseline casuale}      & --- & --- & $\sim$0,50 \\
\bottomrule
\end{tabular}
\end{table}

\paragraph{Preservazione della struttura.}
Al 10\% di contaminazione ($\alpha = 0{,}1$), l'autospettro
dell'adattatore sorgente \`e perfettamente preservato: similarit\`a del
coseno $= 1{,}0000$ per tutte e tre le coppie. Ci\`o vale
simmetricamente ($\alpha = 0{,}9$). Piccole quantit\`a di informazione
inter-task possono essere iniettate nella matrice ponte senza
disturbare la struttura spettrale del task dominante. Anche a
$\alpha = 0{,}5$ (mescolamento paritario), il coseno dell'autospettro
rimane $> 0{,}999$ per tutte le coppie di task. Al contrario, la
baseline casuale produce un coseno $\approx 0{,}5$ a $\alpha = 0{,}5$:
il rumore casuale distrugge la struttura spettrale, mentre
l'interpolazione inter-task la preserva. La fusione delle matrici ponte
\`e strutturata, non rumore.

\paragraph{Sensibilit\`a alla fusione per modulo.}
La sensibilit\`a alla fusione (variazione normalizzata del Fiedler da
$\alpha = 0{,}0$ a $0{,}5$) varia per tipo di modulo. Le proiezioni Q
sono le pi\`u sensibili alla fusione in 2 coppie su 3
(alpaca--codice: 0,429; codice--matematica: 0,160), coerentemente con il
loro status di modulo pi\`u discriminativo per il task
(\S\ref{sec:exp-fingerprint}).
La correlazione cross-fase tra la distanza del task nella Fase~1A e la
sensibilit\`a alla fusione nella Fase~2A \`e $r = 0{,}921$ tra i
quattro tipi di modulo: i moduli che portano la maggiore identit\`a del
task sono i pi\`u vulnerabili all'interferenza inter-task. Le matrici
ponte K-proj e V-proj sono robuste alla fusione, suggerendo una
strategia di composizione pratica: fondere liberamente le matrici ponte
K/V, interpolare con cautela le matrici ponte O, e mantenere matrici
ponte Q specifiche per task.

% ---------------------------------------------------------------------------
\subsection{Diagnostica del Sovradattamento}
\label{sec:exp-overfit}

Addestriamo RhombiLoRA su una
porzione deliberatamente ridotta (500 addestramento / 500 validazione da
Alpaca-cleaned) per 10.000 passi con checkpoint ogni 100 passi,
tracciando il valore di Fiedler e la deviazione insieme al divario di
perdita addestramento--validazione su una RunPod RTX 4090.

La perdita di addestramento scende a 0,012 mentre la perdita di
validazione si stabilizza a 1,51 al passo 10K (divario $= +1{,}50$),
confermando un grave sovradattamento. Entrambe le propriet\`a spettrali
della matrice ponte correlano fortemente con il divario
addestramento--validazione: deviazione ($r = 0{,}888$,
$p = 7{,}3 \times 10^{-35}$) e valore di Fiedler ($r = 0{,}825$,
$p = 5{,}6 \times 10^{-26}$). Entrambi superano la soglia $r > 0{,}7$.

Una transizione di fase nella connettivit\`a della matrice ponte si
verifica al passo 400: la variazione massima del Fiedler tra un passo
e l'altro ($0{,}019$) supera la variazione mediana
($2{,}3 \times 10^{-5}$) di $807\times$, coincidendo con la rapida
discesa iniziale della perdita di addestramento. Dopo questa
riconfigurazione iniziale, le propriet\`a spettrali della matrice ponte
evolvono in modo monotono parallelamente al crescente divario di
generalizzazione. La matrice ponte non predice l'insorgenza del
sovradattamento \emph{in anticipo}---traccia lo stesso segnale di
gradiente che guida il divario di perdita---ma fornisce un riassunto
compatto e privo di inferenza dello stato dell'adattatore che potrebbe
integrare il monitoraggio basato sulla perdita nella pratica.

% ===================================================================
% Section 5 — Lavori Correlati
% ===================================================================

\section{Lavori Correlati}
\label{sec:related}

\paragraph{Varianti di LoRA.}
LoRA~\cite{hu2022lora} ha stabilito la fattorizzazione $B \cdot A$ per
il fine-tuning efficiente in termini di parametri. I lavori successivi
hanno modificato questa fattorizzazione in diverse direzioni.
AdaLoRA~\cite{zhang2023adalora} regola dinamicamente l'allocazione del
rango tra i livelli tramite punteggio di importanza.
LoRA-XS~\cite{banaei2024loraxs} inserisce una matrice $r \times r$ tra
le proiezioni $A$ e $B$ congelate, ottenendo una riduzione estrema dei
parametri apprendendo solo l'accoppiamento intermedio.
MELoRA~\cite{ren2024melora} partiziona le dimensioni di rango in
mini-LoRA indipendenti, equivalenti a una matrice ponte a blocchi
diagonali con blocchi identit\`a congelati.
DoRA~\cite{liu2024dora} decompone l'aggiornamento in componenti di
magnitudine e direzione. Nessuna di queste varianti utilizza la
rappresentazione intermedia come \emph{diagnostica} del comportamento
dell'adattatore; tutte inquadrano la struttura intermedia come un mezzo
per migliorare l'efficienza dell'adattamento.

\paragraph{Fusione degli adattatori.}
L'aritmetica dei task~\cite{ilharco2023editing} compone gli adattatori
mediante combinazione lineare dei delta dei pesi.
TIES-Merging~\cite{yadav2023ties} elimina i parametri di piccola
magnitudine e risolve i conflitti di segno prima della media.
DARE~\cite{yu2024dare} elimina casualmente i parametri dell'adattatore
prima della fusione. AdaMerging~\cite{yang2024adamerging} apprende
coefficienti di fusione per livello. Questi metodi operano su matrici
dei pesi complete o sui loro fattori SVD senza una rappresentazione
compatta dell'identit\`a del task.
L'interpolazione delle matrici ponte (\S\ref{sec:exp-merging}) fornisce
tale rappresentazione: uno spazio a 36 dimensioni in cui la
compatibilit\`a tra adattatori pu\`o essere valutata prima di
impegnarsi in una costosa fusione a peso pieno.

\paragraph{Analisi degli adattatori.}
I lavori precedenti sulla comprensione delle rappresentazioni LoRA si
sono concentrati sulla struttura nello spazio dei pesi. L'analisi delle
dinamiche di addestramento di LoRA~\cite{biderman2024lora} esamina come
$A$ e $B$ evolvono durante l'addestramento ma non fornisce alcun
riassunto compatto dell'accoppiamento appreso. L'analisi spettrale dei
pesi dell'adattatore~\cite{zhang2023adalora} opera sulle matrici di
proiezione complete $d_\text{out} \times r$.
Le nostre diagnostiche basate sulla matrice ponte comprimono la
struttura di accoppiamento in $C^2 = 36$ parametri per livello
mantenendo informazione sufficiente per la classificazione del task
all'84,5\% di accuratezza.

% ===================================================================
% Section 6 — Discussione
% ===================================================================

\section{Discussione}
\label{sec:discussion}

\subsection{Perch\'e la Matrice Ponte Non Migliora la Perdita}

La matrice ponte aggiunge 36 parametri per livello dell'adattatore. In
una configurazione RhombiLoRA a rango 24 applicata a quattro moduli
target su 28 livelli del trasformatore, ci\`o ammonta a
$28 \times 4 \times 36 = 4.032$ parametri della matrice
ponte---trascurabili rispetto al conteggio totale dei parametri
dell'adattatore
($28 \times 4 \times 2 \times 24 \times d_{\text{model}}$, dove
$d_{\text{model}} = 3.584$ per Qwen2.5-7B, producendo circa 19,3M di
parametri dell'adattatore). La matrice ponte contribuisce lo 0,02\%
della capacit\`a totale dell'adattatore.

A questa scala, il panorama di ottimizzazione non beneficia dei gradi di
libert\`a aggiuntivi della matrice ponte. La matrice di accoppiamento
$6 \times 6$ si trova tra due matrici di proiezione molto pi\`u grandi
$A$ e $B$; qualsiasi trasformazione lineare $M$ pu\`o essere assorbita
nel prodotto $BMA$ ridefinendo $B' = BM$. L'ottimizzatore non ha alcun
segnale di gradiente che favorisca la distribuzione della
rappresentazione tra $B$, $M$ e $A$ piuttosto che l'assorbimento di $M$
interamente in $B$. Ci\`o \`e coerente con il risultato del nostro
Esperimento~1: RhombiLoRA con matrice ponte congelata corrisponde alla
perdita del LoRA standard entro $\Delta < 0{,}0002$, e le varianti con
matrice ponte addestrabile non migliorano. Raddoppiare il rango da 24 a
48 non fornisce parimenti alcun beneficio---il collo di bottiglia \`e
il segnale di addestramento, non la capacit\`a dell'adattatore.

\subsection{Perch\'e la Matrice Ponte Cattura Struttura Diagnostica}

Nonostante non contribuisca alla perdita, la matrice ponte cattura
struttura discriminativa per il task poich\'e occupa un collo di
bottiglia rappresentazionale. Tutta l'informazione che fluisce da $A$ a
$B$ passa attraverso la matrice ponte $6 \times 6$ $M$. Durante
l'addestramento, i gradienti che modellano $M$ riflettono il modo in
cui il panorama di perdita del task richiede l'interazione tra i sei
canali di rango. Se un task richiede un ampio mescolamento tra canali
(come nel caso del seguito di istruzioni generali), $M$ sviluppa voci
fuori diagonale; se il task \`e pi\`u vincolato (come nel caso della
generazione di codice), $M$ rimane pi\`u vicino alla diagonale.

Questa \`e la distinzione chiave tra \emph{capacit\`a} e
\emph{diagnostica}. La matrice ponte non ha bisogno di migliorare il
potere rappresentazionale dell'adattatore per essere informativa su ci\`o
che l'adattatore ha appreso. \`E un riassunto a 36 parametri del
pattern di accoppiamento che $A$ e $B$ hanno scoperto
congiuntamente---leggibile senza eseguire l'inferenza, senza valutazione
su dati di riserva e senza accesso all'insieme di addestramento. La
magnitudine dell'accoppiamento (valore di Fiedler) differenzia i task:
gli adattatori alpaca sviluppano un accoppiamento tra canali $2{,}2\times$
superiore rispetto agli adattatori di codice (Fiedler 0,040 contro
0,018). La struttura dell'accoppiamento (forma dell'autospettro) \`e
conservata: la similarit\`a del coseno tra gli autospettri dei task
supera 0,999. I task modulano l'\emph{intensit\`a} dell'accoppiamento
della matrice ponte preservandone l'\emph{architettura}.

\subsection{Il Risultato della Proiezione Q}

In tutte e tre le fasi sperimentali, le matrici ponte della proiezione
query sono la componente pi\`u informativa. Fase~0 (anatomia
single-task): Q-proj sviluppa il 40\% in pi\`u di accoppiamento
rispetto a V-proj ($F = 6{,}022$, $p = 0{,}0008$). Fase~1A
(fingerprinting del task): Q-proj porta la distanza inter-task pi\`u
elevata (0,162 contro 0,075 per K-proj e 0,072 per V-proj). Fase~2A
(fusione degli adattatori): le matrici ponte Q-proj sono le pi\`u
sensibili all'interpolazione inter-task. Questo ordinamento \`e
coerente tra le fasi (Pearson $r = 0{,}921$ tra la discriminazione del
task nella Fase~1A e la sensibilit\`a alla fusione nella Fase~2A).

Il meccanismo \`e architetturale. La proiezione query calcola
\emph{a cosa prestare attenzione}---il calcolo pi\`u specifico per il
task nella testa di attenzione. Le proiezioni chiave calcolano
l'indicizzazione; le proiezioni valore trasportano il contenuto. Le
query codificano la domanda del task sulla struttura dell'input, che
varia massimamente tra i tipi di task. La matrice ponte nella proiezione
Q cattura questa variazione poich\'e si trova nel punto di massima
divergenza dipendente dal task nel flusso di informazione.

Un risultato secondario inatteso: le proiezioni di uscita superano le
proiezioni query negli ultimi cinque livelli del trasformatore (O-proj
del Livello~24 raggiunge la singola distanza inter-task pi\`u elevata
di 0,3188 su tutte le 112 posizioni degli adattatori). Le proiezioni
di uscita dei livelli tardivi riformattano i risultati dell'attenzione
per la predizione a valle, e questa riformattazione diventa fortemente
specifica per il task man mano che il modello si avvicina al suo output.
Il gradiente di profondit\`a \`e significativo (Spearman
$\rho = 0{,}701$, $p < 0{,}0001$), con un intervallo dinamico di
$10{,}3\times$ tra le posizioni di matrice ponte pi\`u e meno
discriminative.

\subsection{Implicazioni Pratiche}

Un SVM leave-one-out addestrato sulle 28 matrici ponte della proiezione
Q di un singolo adattatore classifica il tipo di task con un'accuratezza
dell'84,5\% su tre task (caso: 33,3\%). Questa impronta costa 1.008
parametri ($28 \times 36$)---pi\`u piccola di un singolo vettore di
stato nascosto in Qwen2.5-7B (4.096 dimensioni). Includendo le 28
matrici ponte della proiezione O (2.016 parametri totali) si ottiene
l'83,3\% di accuratezza con prestazioni migliorate sulla classe pi\`u
difficile (matematica: 66,1\% contro 50,0\% con la sola Q-proj).

Per la composizione degli adattatori, l'interpolazione a livello di
matrice ponte preserva la struttura dell'autospettro a tutti i rapporti
di mescolamento (coseno $> 0{,}999$), ma la traiettoria del Fiedler \`e
non lineare per coppie di task strutturalmente dissimili
($R^2 = 0{,}735$ per alpaca$\leftrightarrow$matematica contro
$R^2 = 0{,}948$ per alpaca$\leftrightarrow$codice).
L'interpolazione codice$\leftrightarrow$matematica esibisce una
diminuzione del Fiedler a $\alpha \approx 0{,}2$, coerente con
l'interferenza distruttiva tra pattern della matrice ponte che
differiscono in struttura piuttosto che nella sola magnitudine. Al 10\%
di contaminazione inter-task ($\alpha = 0{,}1$), tutte le coppie
preservano la struttura del task dominante con coseno 1,0000. Ci\`o
fornisce un rapporto di mescolamento sicuro empirico---conservativo, ma
fondato.

Questi risultati suggeriscono un'architettura pratica per la gestione
degli adattatori: memorizzare solo le 28 matrici ponte della proiezione
Q per task (1.008 parametri) come impronta del task; utilizzarle per
indirizzare le richieste in arrivo verso l'adattatore corretto senza
inferenza. Le matrici ponte K-proj e V-proj sono strutturalmente
conservate tra i task e possono essere condivise. Le matrici ponte
Q-proj richiedono un trattamento specifico per task durante la
composizione.

\subsection{Limitazioni}
\label{sec:limitations}

\textbf{Singola famiglia di modelli.} Tutti gli esperimenti utilizzano
Qwen2.5-7B-Instruct. La dominanza della proiezione Q e il gradiente di
profondit\`a potrebbero riflettere la struttura di attenzione di questa
architettura piuttosto che una propriet\`a universale degli adattatori
per trasformatori. L'estensione a LLaMA, Mistral o architetture non
decoder \`e necessaria prima di rivendicare la generalit\`a.

\textbf{Tre task.} La classificazione a tre vie all'84,5\% \`e
incoraggiante ma limitata. Con sole tre categorie di task (generale,
codice, matematica), il confine decisionale \`e a bassa dimensionalit\`a.
Il passaggio a 10--50 task verificherebbe se le impronte delle matrici
ponte rimangono discriminative man mano che il problema di
classificazione diventa pi\`u difficile, o se task strutturalmente simili
(ad es.\ generazione di codice vs.\ debug di codice) producono matrici
ponte indistinguibili.

\textbf{Rango fisso della matrice ponte.} La matrice ponte \`e
$6 \times 6$ in tutti gli esperimenti, motivata dalle sei coppie di
facce del dodecaedro rombico. Ranghi pi\`u elevati della matrice ponte
($8 \times 8$, $12 \times 12$) potrebbero catturare una struttura di
accoppiamento pi\`u fine al costo di un maggiore numero di parametri.
Ranghi inferiori ($3 \times 3$, la baseline cubica) mostrano un
accoppiamento $4{,}6\times$ inferiore nell'Esperimento~1, ma se
mantengano la capacit\`a di fingerprinting non \`e stato verificato.

\textbf{Durate di addestramento non uniformi.} Alpaca \`e stato
addestrato per 10K passi, codice e matematica per 2K ciascuno. La
relazione monotona tra durata dell'addestramento e distanza della
matrice ponte dall'inizializzazione (alpaca: 0,199, matematica: 0,069,
codice: 0,032) implica che parte della separazione inter-task potrebbe
riflettere la durata dell'addestramento piuttosto che la struttura del
task. Il controllo al passo 0 (tutte le distanze esattamente 0,000000)
conferma che le impronte emergono dall'addestramento, ma un confronto
completamente controllato addestrerebbe tutti i task per lo stesso
numero di passi.

\textbf{Nessuna valutazione generativa.} La fusione delle matrici ponte
\`e stata valutata strutturalmente (coseno dell'autospettro, traiettorie
del Fiedler) ma non generativamente. Se lo scambio di matrici ponte tra
adattatori produca un output utilizzabile---se la matrice ponte a 36
parametri porti sufficiente informazione sul task per redirigere una
proiezione $A/B$ fissa---rimane non validato. Questo \`e l'esperimento
critico mancante per l'affermazione sulla composizione degli adattatori.

\textbf{Metodologia di classificazione.} L'SVM leave-one-out su 84
campioni (28 per task) \`e una valutazione a piccolo campione. Il
classificatore utilizza vettori di caratteristiche a 36 dimensioni
($6 \times 6$) appiattiti senza ingegneria delle caratteristiche
informata dal dominio. Approcci pi\`u sofisticati (caratteristiche
spettrali, metriche grafo-teoriche, rappresentazioni
GL($r$)-equivarianti~\cite{putterman2024learning}) potrebbero migliorare
l'accuratezza ma complicherebbero l'argomento di interpretabilit\`a.

\subsection{Connessione con la Topologia Reticolare}
\label{sec:lattice-connection}

La matrice ponte \`e il punto in cui la geometria reticolare degli
Articoli~1--2~\cite{rhombilora_paper1, rhombilora_paper2} incontra
l'architettura neurale. Le sei coppie di facce del dodecaedro rombico
definiscono sei coppie di canali nella matrice ponte; la connettivit\`a
del reticolo FCC determina la topologia di inizializzazione.
L'Articolo~1 ha stabilito che i reticoli FCC raggiungono una
connettivit\`a algebrica $2{,}4\times$ superiore rispetto ai reticoli
cubici a parit\`a di numero di nodi. L'Articolo~2 ha mostrato che
questo vantaggio si amplifica a $6{,}1\times$ sotto pesi degli archi
strutturati. La matrice ponte \`e il sito in cui queste propriet\`a
reticolari entrano nella rete neurale.

La geometria reticolare fornisce due risorse distinte.
\textbf{Connettivit\`a}---il vantaggio algebrico di pi\`u percorsi di
accoppiamento---emerge naturalmente dall'addestramento non
supervisionato. Sei canali con inizializzazione a topologia FCC
sviluppano una connettivit\`a algebrica $4{,}6\times$ superiore rispetto
a tre canali con inizializzazione cubica a parit\`a di budget di
parametri (\S\ref{sec:exp-anatomy-1b}).
L'ottimizzatore scopre indipendentemente che la disposizione FCC dei
percorsi di accoppiamento \`e utile, senza alcuna supervisione sulla
semantica dei canali. Questa \`e la dimensione della capacit\`a: la
topologia determina \emph{quanto} accoppiamento la matrice ponte
sviluppa.

\textbf{Direzione}---l'associazione di canali specifici con direzioni
specifiche del reticolo---non pu\`o emergere dai soli dati. Due
risultati nulli indipendenti lo confermano: dati geometrici strutturati
producono Co/Cross $= 1{,}002$ (\S\ref{sec:exp-direction}); un tasso di
apprendimento della matrice ponte $10\times$ pi\`u elevato produce
Co/Cross $\approx 1{,}0$ (\S\ref{sec:exp-bridge-lr}). Tuttavia, 500
passi di supervisione contrastiva installano una preferenza direzionale
(Co/Cross $= 2.660\times$) che persiste come memoria geometrica stabile:
dopo 6.000 passi di addestramento non supervisionato, il rapporto rimane
a $16\times$---senza mai avvicinarsi alla baseline nulla
(\S\ref{sec:exp-contrastive}). Durante lo stesso periodo non
supervisionato, la connettivit\`a di Fiedler cresce di $3{,}7\times$
indipendentemente. La matrice ponte sviluppa simultaneamente pi\`u
accoppiamento complessivo \emph{e} mantiene la struttura direzionale che
le \`e stata insegnata.

Questa separazione tra le due risorse chiarisce la relazione tra
topologia reticolare e adattamento neurale. Il reticolo FCC fornisce un
\emph{sistema di coordinate geometriche indirizzabile}: alle sue sei
coppie di direzioni pu\`o essere assegnato un significato semantico
attraverso la supervisione contrastiva, e quel significato persiste
attraverso l'addestramento successivo. La connettivit\`a \`e la risorsa
gratuita---qualsiasi topologia la fornisce, ma l'FCC ne fornisce
$4{,}6\times$ di pi\`u. La direzione \`e la risorsa costosa---deve
essere esplicitamente insegnata. Una volta insegnata, la geometria
reticolare garantisce che abbia un luogo stabile in cui risiedere.

% ===================================================================
\section{Conclusione}
\label{sec:conclusion}

Abbiamo introdotto una matrice di accoppiamento addestrabile
$6 \times 6$---la \emph{matrice ponte}---tra le proiezioni $A$ e $B$ di
un adattatore LoRA, aggiungendo 36 parametri per livello. La matrice
ponte non migliora la perdita di fine-tuning: a rango 24, i suoi
parametri sono una frazione trascurabile della capacit\`a totale
dell'adattatore, e il panorama di ottimizzazione non beneficia dei gradi
di libert\`a aggiuntivi. Non si tratta di un adattatore migliore. \`E
una nuova lente sul comportamento dell'adattatore.

La matrice ponte rivela ci\`o che LoRA apprende, in 36 parametri per
livello. I task che richiedono un ampio mescolamento tra canali (seguito
di istruzioni generali) producono matrici ponte con una connettivit\`a
algebrica $2{,}2\times$ superiore rispetto ai task con struttura
vincolata (generazione di codice). Un SVM leave-one-out addestrato su
28 matrici ponte della proiezione query---1.008 parametri
totali---classifica il tipo di task con un'accuratezza dell'84,5\% senza
inferenza n\'e valutazione su dati di riserva. L'impronta emerge
interamente dall'addestramento: al passo~0, tutte le matrici ponte sono
identiche; al passo~1.500, la classificazione binaria tra i due task
pi\`u simili raggiunge il 93\%.
La proiezione query \`e la posizione ottimale per l'impronta poich\'e
codifica l'aspetto pi\`u specifico per il task dell'attenzione: cosa
cercare.

La composizione degli adattatori a livello di matrice ponte preserva la
struttura dell'autospettro a tutti i rapporti di mescolamento
(coseno $> 0{,}999$), ma le traiettorie non lineari del Fiedler nelle
coppie di task strutturalmente dissimili rivelano interferenza
distruttiva invisibile alle diagnostiche di fusione a peso pieno. La
matrice ponte decompone il comportamento dell'adattatore in una
componente di \emph{proiezione} ($A$ e $B$, milioni di parametri) e una
componente di \emph{accoppiamento} ($M$, 36 parametri)---e la sola
componente di accoppiamento porta struttura sufficiente per identificare
il task, rilevare l'interferenza durante la fusione e differenziare i
componenti della testa di attenzione per funzione.

La matrice ponte esibisce inoltre \emph{memoria geometrica}. Il
pre-addestramento contrastivo installa una preferenza direzionale nella
struttura di accoppiamento della matrice ponte (rapporto
co-planare/cross-planare $> 2.600\times$ dopo 500 passi
supervisionati). Quando la supervisione viene rimossa, il segnale
direzionale decade ma non raggiunge mai la baseline nulla: dopo 6.000
passi non supervisionati, il rapporto si stabilizza intorno a
$16\times$---mentre due esperimenti di controllo indipendenti confermano
che l'addestramento non supervisionato da solo produce un rapporto di
$\approx 1{,}0$.
Durante lo stesso periodo non supervisionato, la connettivit\`a
complessiva della matrice ponte cresce di $3{,}7\times$
indipendentemente: l'ottimizzatore scopre autonomamente l'utilit\`a
dell'accoppiamento tra canali, ma non pu\`o scoprire \emph{quali}
canali accoppiare preferenzialmente. La matrice ponte separa due risorse
fornite dalla topologia reticolare: connettivit\`a (gratuita, emerge
dall'addestramento) e direzione (costosa, richiede insegnamento,
persiste una volta insegnata).

Il presente studio \`e limitato a una singola famiglia di modelli, tre
task e valutazione strutturale senza validazione generativa. Il
passaggio a pi\`u task verificher\`a se le impronte delle matrici ponte
rimangono discriminative in problemi di classificazione a dimensionalit\`a
pi\`u elevata. Ranghi pi\`u elevati della matrice ponte potrebbero
catturare una struttura di accoppiamento pi\`u fine. Il risultato sulla
memoria geometrica apre una nuova direzione: se un breve riscaldamento
contrastivo pu\`o installare una struttura geometrica persistente, la
matrice ponte potrebbe servire non solo come diagnostica passiva ma come
sito per l'iniezione di bias induttivi che modellano l'addestramento non
supervisionato successivo. La direzione pratica pi\`u immediata \`e
l'instradamento degli adattatori basato sulla matrice ponte: data una
libreria di adattatori specifici per task, selezionare quello corretto
confrontando l'impronta della proiezione Q a 1.008 parametri di un
adattatore interrogativo con le impronte memorizzate, senza eseguire
alcuno degli adattatori al momento dell'inferenza.

\paragraph{Disponibilit\`a.}
Codice sorgente: \url{https://github.com/promptcrafted/rhombic}.\\
Pacchetto: \url{https://pypi.org/project/rhombic/}.\\
Matrici ponte e script di analisi: inclusi nel repository sotto
\texttt{results/} e \texttt{scripts/}.

\section*{Ringraziamenti}
Questo lavoro \`e stato condotto interamente su hardware GPU consumer e
cloud: NVIDIA RTX 6000 Ada 48\,GB (locale), RTX 4090 16\,GB e RTX 3090
24\,GB (RunPod) per esperimenti su piccola scala e contrastivi. La
libreria \texttt{rhombic} e tutto il codice sperimentale sono
disponibili presso \url{https://github.com/promptcrafted/rhombic} sotto
licenza MPL-2.0. Le matrici ponte per tutti gli esperimenti sono incluse
come array NumPy per la piena riproducibilit\`a.

\bibliographystyle{plainnat}
\bibliography{refs-paper3}

\end{document}
